% Auxiliary statements for Proposition 9.2, not the complete scanner estimate.
\section{Charge selection with an arbitrary fixed weight}
\label{sec:arealaw-source-charge-selection}

The construction below is an auxiliary step in the proof of
\cite[Proposition~9.2]{OpenAI2026AreaLaw}. It retains the physical
state and label sequence while allowing any positive total weight and
any fixed interpolation interval contained in $(0,1)$.

\begin{theorem}[Nonvanishing and the eigen-equation in augmented coordinates]
    \label{thm:arealaw-augmented-prevector-properties}
    \lean{TNLean.PEPS.AreaLaw.Scan.augmented_prevector_properties_of_global}
    \leanok
    \uses{def:al_scan_auxiliary_coordinates}
    Let $V$ be finite, let $q,d,k\in\mathbb N$, let $\Omega$ be a complex
    vector on the physical configurations, and let $\lambda$ be a label
    of the symmetric group on $k$ copies. Let $v$ be the corresponding
    physical Schmidt--Bell prevector in the global coordinates, and let
    $\mathrm{pre}=Uv$ be its expression in the augmented site coordinates,
    where $U$ is the coordinate isometry. For a physical matrix $H$, write
    $H_{\rm glob}$ and $H_{\rm aug}$ for its replica mean, with identities
    on the two auxiliary registers, in these two coordinate systems.
    If $v\ne0$ and $H_{\rm glob}v=e v$ for a real $e$, then
    \begin{align}
        \mathrm{pre}            & \ne0,\notag                    \\
        \mathrm{pre}            & \in\mathcal H_{\rm sym},\notag \\
        H_{\rm aug}\mathrm{pre} & =e\,\mathrm{pre}.\notag
    \end{align}
    Here $\mathcal H_{\rm sym}$ is the subspace invariant under simultaneous
    permutation of the copies. This coordinate statement is used in
    \cite[Section~8, lines~130--153 and~421--438]{OpenAI2026AreaLaw}.
\end{theorem}
\begin{proof}
    \leanok
    \uses{thm:al_augmented_prevector_symmetry,thm:al_scan_global_replica_energy}
    The coordinate bijection preserves nonvanishing. The augmented
    prevector is fixed by the simultaneous symmetry projection, so it
    lies in its range. The coordinate identity
    $H_{\rm aug}=UH_{\rm glob}U^{-1}$ gives
    $H_{\rm aug}\mathrm{pre}=UH_{\rm glob}v=e\,\mathrm{pre}$.
\end{proof}

\begin{theorem}[Selected density from the rough norm difference]
    \label{thm:arealaw-density-from-rough-endpoint}
    \lean{TNLean.PEPS.AreaLaw.Scan.selected_density_of_rough_endpoint}
    \leanok
    \uses{def:arealaw-rough-endpoint-remainder}
    Let $k,n,m,D,K$ be positive integers, let $d\in\mathbb N$, and let
    $a,C_1,w>0$. Put $W=aK$ and $\rho=ka/(2(C_1+1)nD)$.
    For real $A,B,T,b,c,C_i,C_b,\mathcal Q$, suppose
    \begin{align}
        \frac A k  & \leq B+R_k^{\rm end}(d,c,K,C_i),\notag  \\
        b          & \leq C_b\log(k+2),\notag                \\
        \mathcal Q & \leq\frac{A+2nm(kT+b)}{\rho nmw}.\notag
    \end{align}
    Then
    \begin{align}
        \frac{\mathcal Q}{KnD}
         & \leq\frac{2(C_1+1)}{Wnmw}
        \left[B+2nmT+R_k^{\rm end}(d,c,K,C_i+2nmC_b)\right].\notag
    \end{align}
    This is the scalar normalization of the selected charge in
    \cite[Section~9, lines~457--490]{OpenAI2026AreaLaw}.
\end{theorem}
\begin{proof}
    \leanok
    \uses{thm:arealaw-rough-selection-numerator,thm:arealaw-selected-density-normalization}
    The rough norm difference and logarithmic error bounds give
    $(A+2nm(kT+b))/k\leq B+2nmT+R_k^{\rm end}(d,c,K,C_i+2nmC_b)$.
    Normalize the selected inequality by $KnD$, cancel the positive
    factors in $\rho$, and use $aK=W$.
\end{proof}

\begin{theorem}[One physical state and label sequence for charge selection]
    \label{thm:arealaw-source-charge-selection}
    \lean{
        TNLean.PEPS.AreaLaw.Scan.CollarScan.exists_sourceTarget_selected_actualCharge_density_energy}
    \leanok
    \uses{def:area_law_hamiltonian,def:al_scan_actual_histories,
        def:al_scan_actual_energy_terms}
    Fix an integer $q\geq1$, an interaction range $R\in\mathbb N$, $J\geq0$,
    $\Delta>0$, $C_T\geq0$, $C_1>0$, an integer anchor-multiplicity
    bound $m_*\geq0$ with $3m_*\leq C_1$, and a real exponent $P$.
    There are $C_{\rm tr},C_{\rm target},c_t>0$ and real
    $C_{\rm ent},e_{\rm ent}$ with the following property. For every
    scanner exponent tuple, there are $C\geq1$, $C_l\geq0$ and
    $N\geq2$ before all the physical data and entropy inputs below.

    Let $\Lambda,T\subset\mathbb Z^2$ be finite, with $T\ne\varnothing$,
    and let $H$ be a range-$R$ Hamiltonian of interaction
    bound $J$, with unit ground vector $\Omega$ and projector gap
    $\Delta$. Use the actual induced graph, anchors belonging to their
    interaction supports, and depth $d(x)=\operatorname{dist}_\infty(x,T)$.
    For the scan on $A\subseteq\Lambda$, assume $n\geq N$ and the scales
    \begin{align}
        r_0 & =\lceil C_{\rm tr}(\log n)^2\rceil,\notag     \\
        L   & =\lfloor n^{1-\ell}\rfloor,\notag             \\
        m   & =\lfloor n^\mu\rfloor,\notag                  \\
        D   & =\lceil n^\kappa\rceil,\notag                 \\
        K   & =\left\lfloor\frac{L}{8m}\right\rfloor.\notag
    \end{align}
    The exponent conditions are $0<\ell<1$, $0<\kappa<\mu<1-\ell$,
    $\nu>0$ and $0<e_{\rm in}<1$ with $\kappa\leq e_{\rm in}$.
    The charge constant is $C_1$. Put $M=\lceil C_1nD\rceil$.
    Suppose $|T|\leq C_Tn^2$, every
    ambient dilation layer through depth $L$ has at most $n$ sites,
    and at most $m_*$ interaction supports have any given anchor.
    Every point of $T$ is at distance greater than $2L+10r_0$ from
    every endpoint of a cut edge between $A$ and its complement.
    Also suppose $2r_0\leq D$, $1\leq D$ and $4D\leq m$.
    For arbitrary real $C_e,\eta$, assume
    $S_\Omega(\{x\in A:0<d(x)\leq j\})\leq C_e(nL^\eta+n)$
    for every integer $0\leq j\leq L$.

    Put $X=A\cap T$ and $S_0=A\cap T_L$, where $T_L$ is the ambient
    dilation of $T$. Let $g>0$ be the normalized amplified gap in
    Theorem~\ref{thm:al_scan_target_typical_window}, and set
    \begin{align}
        \varepsilon_t & =\min\{n^{-1000},g/4\},\notag \\
        w             & =n^{3/5},\notag               \\
        \delta_t      & =n^{-100}.\notag
    \end{align}
    Then $m,M>0$, $K\geq8$, $8Km\leq L$, $|S_0|\leq(C_T+1)n^2$
    and $X\subseteq S_0$. There is a unit ground vector
    $\widetilde\Omega$ of the truncated Hamiltonian $\widetilde H$,
    with energy $e_t$ and projector gap $g/2$, such that
    \begin{align}
        0\leq e_t    & \leq\varepsilon_t,\notag                 \\
        \operatorname{dist}_{\rm tr}
        (|\widetilde\Omega\rangle\langle\widetilde\Omega|,
        |\Omega\rangle\langle\Omega|)
                     & \leq\sqrt{2\varepsilon_t/g},\notag       \\
        \mathcal B_X & \leq C_{\rm target}n(\log n)^{12}.\notag
    \end{align}
    Write $p_i$ for the target-marginal eigenvalues and
    $S_X=S_{\widetilde\Omega}(X)$. Their surprisal tail outside width
    $w$ has mass at most $\delta_t$. The typical set
    $E=\{i:p_i>0,\ |{-\log p_i}-S_X|\leq w\}$ has cardinality $d>0$
    and mass $z$ with $0<1-\delta_t\leq z\leq1$; every $i\in E$
    obeys $\exp(-S_X-w)\leq p_i\leq\exp(-S_X+w)$.

    Let $\psi$ be the normalized Schmidt vector restricted to $E$, and
    let $S'$ be its auxiliary entropy. One sequence $\lambda_k$ and
    its physical Schmidt--Bell prevectors $\mathrm{pre}_k$ satisfy
    $\mathrm{pre}_k\ne0$ and $\|\mathrm{pre}_k\|\leq1$ for all $k>0$,
    and $\log\dim\lambda_k-kS'=o(k)$. For all sufficiently large $k$,
    \begin{align}
        \|(I\otimes P_{\lambda_k})\psi^{\otimes k}\|^2
         & \geq\frac1{2(k+1)^{d^2}}.\notag
    \end{align}
    The state $\widetilde\Omega$, the set $E$, the compressed vector $\psi$
    and the sequences $(\lambda_k)$ and $(\mathrm{pre}_k)$ are fixed before $W$.
    For any $W>0$, put $a=W/K$ and
    $\ell_{\dim}=1+(2r_0+1)^2\log q$, and suppose
    $a\ell_{\dim}\leq c_t$. Define
    \begin{align}
        B_{\rm sh}  & =C_e(nL^\eta+n)+n^{-247},\notag                   \\
        B_{\rm rough}
                    & =2\log d-\log z+
        W\left(\frac{B_{\rm sh}}z+1+3nD\log q\right),\notag             \\
        T_{\rm ent} & =C_{\rm ent}aK\,a^{1/4}\ell_{\dim}^{e_{\rm ent}}.
        \label{eq:arealaw-charge-budgets}
    \end{align}
    There are nonnegative $C_b,c_U,c_R,C_i$, a sequence
    $0\leq b_k\leq C_b\log(k+2)$, and $r_k\geq0$ tending to zero.
    Write
    \begin{align}
        R_k^{\rm Bell}(d,c,K)
         & =\frac{\log4+(d^2+cK)\log(k+2)}k,\notag \\
        R_k^{\rm end}
         & =R_k^{\rm Bell}(d,c_R,K)
        +(C_i+2nmC_b)\frac{\log(k+2)}k.\notag
    \end{align}
    The remainders $R_k^{\rm Bell}$ and $R_k^{\rm end}$ tend to zero. For every fixed
    $0<p_0<p_1<1$ and all sufficiently large $k$, there is a charge
    round $s_k<nm$ and $p_k\in[p_0,p_1]$ with the following properties.
    Let $\mathcal Q_k$ be the actual entropy defect at this charge, and
    let $\mathcal M_k$ be its root metric at $p_k$. Define
    \begin{align}
        N_k^2 & =\|\mathcal M_k^{-1/4}\mathrm{pre}_k\|^2,\notag \\
        v_k   & =\mathcal M_k^{-1/4}\mathrm{pre}_k/N_k,\notag   \\
        \mu_k & =\operatorname{Re}\langle v_k,
        \overline{\widetilde H}_k v_k\rangle.\notag
    \end{align}
    Here $N_k$ is the positive square root of $N_k^2$, whose positivity
    is part of the conclusion, and $\overline{\widetilde H}_k$ is the
    mean truncated physical Hamiltonian on the replicas, extended by the
    identity on both auxiliary registers. Put
    \begin{align}
        \operatorname{cap}_k
         & =CW^2n^\ell D^2(\log n)^{C_l}
        \left[(1-p_k)\left(\frac{\mathcal Q_k}{KnD}\right)^{1/8}
            +(1-p_k)n^{-P}+p_k\right].\notag
    \end{align}
    At this single point,
    \begin{align}
        \frac{\mathcal Q_k}{KnD}
              & \leq\frac{2(C_1+1)}{Wnm(p_1-p_0)}
        \left(B_{\rm rough}+2nmT_{\rm ent}+R_k^{\rm end}\right),\notag \\
        \frac{\mu_k-e_t}{g/2}
              & \leq\frac{e_t+\operatorname{cap}_k+r_k}{g/2},\notag    \\
        N_k^2 & >0,\notag                                              \\
        -\frac{\log N_k^2}{k}
              & \leq B_{\rm rough}+R_k^{\rm Bell}(d,c_U,K).
        \label{eq:arealaw-charge-selection}
    \end{align}
    The remainder sequences and their coefficients are chosen before the
    interpolation interval. They may depend on the fixed physical instance
    and weight; the selected round and parameter may depend on $k$.
\end{theorem}
\begin{proof}
    \leanok
    \uses{thm:al_scan_target_typical_window,thm:al_global_schmidt_bell_prevector,
        thm:al_positive_prefix_entropy_transfer,thm:al_scan_actual_rough_upper,
        thm:arealaw-physical-endpoint-loss,thm:arealaw-density-from-rough-endpoint,
        thm:arealaw-augmented-prevector-properties,thm:arealaw-actual-selected-energy,
        thm:arealaw-generic-charge-energy,thm:arealaw-generic-selected-energy,
        lem:al_scanner_scale_facts,thm:area_law_band_count_growth,
        thm:arealaw-rough-endpoint-remainder-limit}
    Choose the truncation constants, then the energy coefficient and size
    threshold for the given exponents. The target-window theorem gives
    $\widetilde\Omega,E,z$; the common-label theorem gives
    $\lambda_k$, its projected mass and its dimension asymptotic.
    Entropy continuity transfers every original prefix bound to the same
    truncated vector with the additive error $n^{-247}$.

    Apply the rough upper estimate and the initial-to-final norm estimate
    to this sequence. After fixing the finite physical instance, weight
    and interval, intersect these two events with $k>0$.
    The augmented-coordinate theorem preserves the nonzero symmetric
    prevector and its mean-energy eigen-equation, so the charge selector
    supplies one point with both the entropy-defect and normalized energy
    inequalities. Apply Theorem~\ref{thm:arealaw-density-from-rough-endpoint},
    with the accumulated logarithmic transport error, to obtain the first
    inequality of~(\ref{eq:arealaw-charge-selection}). At this same selected
    point, the generic coefficient bound replaces the transport-error term
    in the normalized energy inequality by $\operatorname{cap}_k$.
    The rough upper estimate supplies the last two conclusions at this charge.
\end{proof}

\begin{theorem}[Moment radii from marginal smallness]
    \label{thm:arealaw-sharp-smallness}
    \lean{TNLean.PEPS.AreaLaw.Scan.sharpSmallness_le_moment_radius}
    \leanok
    Let $B,a,c_0\in\mathbb R$. If $B\geq1$ and
    $a\sqrt B\leq\min\{c_0,1\}/10$, then
    $10a\leq c_0/\sqrt B$ and $10a\leq1$.
    This is the elementary smallness implication used in
    \cite[Proposition~9.2, sharp-comparison clause]{OpenAI2026AreaLaw}.
\end{theorem}
\begin{proof}
    \leanok
    Since $\sqrt B\geq1$, it is positive. Multiplication of the smallness
    inequality by $10$ gives $10a\sqrt B\leq\min\{c_0,1\}$.
    Division by $\sqrt B$ yields the first bound and
    $10a\leq1/\sqrt B\leq1$.
\end{proof}

% Auxiliary density consequence for Proposition 9.2.
\begin{theorem}[Physical selected density for arbitrary scanner exponents]
    \label{thm:arealaw-source-selected-density}
    \lean{
        TNLean.PEPS.AreaLaw.Scan.CollarScan.exists_sourceTarget_selected_actualCharge_source_density}
    \leanok
    \uses{thm:arealaw-source-charge-selection}
    Fix integers $q\geq1$, $R\geq0$, and real $J\geq0$, $\Delta>0$, $C_T\geq0$,
    $C_1>0$, an anchor-multiplicity bound $m_*\in\mathbb N$ with $3m_*\leq C_1$,
    and a real exponent $P$. There are positive truncation constants
    and $0<c\leq1/2$, chosen before the scanner exponents and entropy
    coefficient. For every admissible exponent tuple
    $(\ell,\mu,\kappa,\nu,e_{\rm in})$, there are $C\geq1$ and $C_l\geq0$.
    For every $C_e\geq0$, there are $C_\delta>0$, $C_{l,\delta}\geq0$
    and $N\geq2$, all chosen before the finite physical instance and $W$.

    Use the original gapped Hamiltonian, nonempty target, actual distance
    function, anchor multiplicity, layer-cardinality bound, cut clearance,
    and rounded scales of Theorem~\ref{thm:arealaw-source-charge-selection}.
    Suppose $n\geq N$, and put $X=A\cap T$. For every integer $0\leq j\leq L$,
    assume
    \begin{align}
        S_\Omega(X) & \leq C_en^{1+e_{\rm in}},\notag    \\
        S_\Omega(\{x\in A:0<d(x)\leq j\})
                    & \leq C_e(nL^{e_{\rm in}}+n).\notag
    \end{align}
    Then $2r_0\leq D$, $1\leq D$, and $4D\leq m$. There is one truncated
    vector $\widetilde\Omega$, typical set $E$, positive selected mass $z$,
    and sequence of labels and physical Schmidt--Bell prevectors satisfying
    all the state, projected-mass and dimension conclusions of
    Theorem~\ref{thm:arealaw-source-charge-selection}, before $W$ is chosen.

    For any $W\geq1$, set $a=W/K$, $\varepsilon=n^{-\nu}$, and assume
    $ar_0^2\leq c$. With the same $B_{\rm rough}$ and $T_{\rm ent}$ as in
    (\ref{eq:arealaw-charge-budgets}), with $\eta=e_{\rm in}$, one has
    \begin{align}
        B_{\rm rough}
         & \leq (6C_e+13+6\log q)Wn^{1+e_{\rm in}},\notag \\
        \delta_n
         & :=C_\delta\varepsilon^{-1}
        \left(n^{e_{\rm in}-\mu}+a^{1/4}(\log n)^{C_{l,\delta}}\right).
        \label{eq:arealaw-source-density}
    \end{align}
    There is a nonnegative sequence $\rho_k\to0$ such that, for all
    sufficiently large $k$, the same actual selected charge $s_k<nm$ and
    parameter $p_k\in[\varepsilon/2,\varepsilon]$ satisfy
    \begin{align}
        \frac{\mathcal Q_k}{KnD} & \leq\delta_n+\rho_k.\notag
    \end{align}
    At this same point all four conclusions of
    (\ref{eq:arealaw-charge-selection}) hold, with
    $p_0=\varepsilon/2$ and $p_1=\varepsilon$. In particular, the actual
    energy cap retains its literal density and truncation term, and the
    filtered prevector has positive norm. Every remainder tends to zero
    with the finite instance and $W$ fixed. The exponent conditions do not
    require $\delta_n\to0$. The case $C_e=0$ is included.
    This is the density consequence used in
    \cite[Proposition~9.2, lines~475--490]{OpenAI2026AreaLaw}.
\end{theorem}
\begin{proof}
    \leanok
    \uses{thm:arealaw-source-charge-selection,thm:scanner_rounded_admissibility,
        thm:al_scan_radius_dimension_smallness,thm:al_scan_transport_coefficients_log,
        thm:arealaw-target-entropy-transfer,thm:arealaw-source-rough-budget-power,
        thm:arealaw-selected-density-delta}
    Enlarge the size threshold so that the rounded scales are admissible,
    $r_0\geq1$, $\log n\geq1$, and truncation changes each relevant entropy
    by at most $n^{-247}$. The radius bound gives the transport smallness
    required for the common state and label sequence. Since $\kappa\leq e_{\rm in}$,
    the ceiling estimate gives $D\leq2n^{e_{\rm in}}$. The target entropy,
    selected-mass estimate and logarithmic dimension bound therefore yield
    the asserted rough-budget bound.

    Write $G=C_{\rm ent}\ell_{\dim}^{e_{\rm ent}}$ and choose uniform
    $C_t\geq1$, $C_{l,\delta}\geq0$ with
    $G\leq C_t(\log n)^{C_{l,\delta}}$. If $C_s\geq1$ is the rounded-scale
    constant, take
    \begin{align}
        C_\delta & =4(C_1+1)(6C_e+13+6\log q+2C_t)C_s,\notag   \\
        \gamma   & =\frac{2(C_1+1)}{Wnm(\varepsilon/2)},\notag \\
        \rho_k   & =|\gamma R_k^{\rm end}|.\notag
    \end{align}
    The identity $aK=W$ gives $T_{\rm ent}=GW a^{1/4}$. Apply the scalar
    density estimate to the rough budget at the point already selected by
    Theorem~\ref{thm:arealaw-source-charge-selection}. Its remaining term
    is $\gamma R_k^{\rm end}\leq\rho_k$. The coefficient $\gamma$ is fixed
    before $k\to\infty$, so $\rho_k\to0$. The energy and norm conclusions
    at that point are unchanged.
\end{proof}
